\chapter{Results}
This chapter presents the experimental findings of our evaluation. We detail the performance of the extraction and classification modules, assess the quality of the hybrid retrieval strategy, and analyze the system's computational efficiency and operational costs.

\section*{Chapter Outline}
This chapter is organized as follows:
\begin{itemize}
    \item \textbf{Section 11.1: Extraction Results} reports on the accuracy of deterministic field extraction.
    \item \textbf{Section 11.2: Classification Results} evaluates the semantic mapping of roles and skills.
    \item \textbf{Section 11.3: Retrieval Results} compares the proposed hybrid search against lexical baselines.
    \item \textbf{Section 11.4: Efficiency and Cost} analyzes system throughput, latency, and resource utilization.
\end{itemize}

\section{Extraction Results}
The first stage of evaluation focuses on the performance of the rule-based Layer A for structured field extraction. Table \ref{tab:extraction_results} summarizes the results across the 500-document "Gold Standard" set.

\begin{table}[H]
    \centering
    \caption{Performance of deterministic field extraction (Layer A)}
    \label{tab:extraction_results}
    \begin{tabular}{@{}lccc@{}}
        \toprule
        \textbf{Field} & \textbf{Precision} & \textbf{Recall} & \textbf{F1-Score} \\ \midrule
        Location (City/Country) & 0.94 & 0.88 & 0.91 \\
        Compensation (Numeric) & 0.96 & 0.92 & 0.94 \\
        Dates (Start/End) & 0.89 & 0.85 & 0.87 \\
        Skill Match (Lexical) & 0.91 & 0.76 & 0.83 \\ \bottomrule
    \end{tabular}
\end{table}

The system achieves high precision across all fields, particularly for compensation and locations where patterns are relatively stable. The lower recall for "Skill Match" is expected, as simple lexicon lookups cannot capture skills described through paraphrasing, which are later addressed in Layer B.

\section{Classification Results}
Layer B's primary task is to bridge the semantic gap using small transformer models. We evaluate the mapping of raw job titles to canonical ESCO occupations and the categorization of complex skill mentions.

\begin{table}[H]
    \centering
    \caption{Classification accuracy for semantic mapping (Layer B)}
    \label{tab:classification_results}
    \begin{tabular}{@{}lcc@{}}
        \toprule
        \textbf{Task} & \textbf{Top-1 Accuracy} & \textbf{Top-3 Accuracy} \\ \midrule
        Title to ESCO Occupation & 0.78 & 0.91 \\
        Skill to ESCO Pillar & 0.84 & 0.93 \\
        Seniority Level & 0.88 & - \\ \bottomrule
    \end{tabular}
\end{table}

The Top-3 accuracy of 91\% for title normalization indicates that while the system may occasionally miss the exact canonical label, it consistently identifies the correct professional family. This is a significant improvement over simple keyword-based title matching, which achieved only 54\% Top-1 accuracy in our baseline tests.

\section{Retrieval Results}
The retrieval quality is measured by comparing our proposed Hybrid Retrieval strategy (Baseline 3) against the Keyword-Only baseline (Baseline 1).

\begin{table}[H]
    \centering
    \caption{Retrieval performance comparison}
    \label{tab:retrieval_results}
    \begin{tabular}{@{}lccc@{}}
        \toprule
        \textbf{Method} & \textbf{P@5} & \textbf{MAP} & \textbf{NDCG} \\ \midrule
        Baseline 1: Keyword-Only & 0.62 & 0.58 & 0.64 \\
        Baseline 3: Hybrid (Proposed) & 0.84 & 0.79 & 0.86 \\ \bottomrule
    \end{tabular}
\end{table}

The hybrid approach demonstrates a substantial gain of +22\% in P@5 and +21\% in MAP. Qualitative analysis shows that the hybrid system excels in handling conceptual queries like "Data Analysis roles in Healthcare," where the keyword baseline fails to find relevant posts that use synonymous terms like "Clinical Researcher" or "Bioinformatics Intern."

\section{Efficiency and Cost}
A key goal of this thesis is to demonstrate the feasibility of running a semantic pipeline on commodity hardware. Table \ref{tab:efficiency_results} reports the performance measured on a machine with 16GB RAM and a 4-core CPU (no dedicated GPU).

\begin{table}[H]
    \centering
    \caption{System efficiency and operational costs}
    \label{tab:efficiency_results}
    \begin{tabular}{@{}ll@{}}
        \toprule
        \textbf{Metric} & \textbf{Value} \\ \midrule
        Ingestion Throughput & 12.5 documents/second \\
        Average Search Latency & 45 ms \\
        Memory Usage (DB + Index + Models) & 3.2 GB \\
        Estimated Cost per 10k Documents (Layer B) & \$0 (local) \\
        Estimated Cost per 10k Documents (Layer C - GPT-4o-mini) & \textasciitilde \$15.00 \\ \bottomrule
    \end{tabular}
\end{table}

By utilizing Layer B for 95\% of the workload and reserving Layer C for ambiguous cases, the system reduces API costs by over 90\% compared to an LLM-only approach, while maintaining a high ingestion throughput suitable for real-time web scraping updates.
